\section{Original homogeneous forms as actual original projective line sections}

\begin{lemma}[The whole original homogeneous fractions have the original scalar transitions]
\label{lem:fc-original-homogeneous-form-coefficients}
\leanok
\lean{EnvelopeProperty.ProjectiveHomogeneousSections.fraction, EnvelopeProperty.ProjectiveHomogeneousSections.coefficient, EnvelopeProperty.ProjectiveHomogeneousSections.fraction_overlap, EnvelopeProperty.ProjectiveHomogeneousSections.coefficient_X, EnvelopeProperty.ProjectiveHomogeneousSections.coefficient_overlap}
For every homogeneous polynomial $F$ of degree $d$ over an arbitrary coefficient
ring, $F/X_i^d$ is an actual entire original coordinate-chart section.  On the
entire original $i,k$ overlap these original sections satisfy
$F/X_i^d=(X_k/X_i)^d(F/X_k^d)$, with the SAME original coordinate-transition unit.
For a degree-one coordinate form the fraction is the literal original
coordinate-hyperplane local equation, including the self-coordinate case.
\end{lemma}
\begin{proof}
Construct the original fractions in the actual homogeneous localizations.
The two genuine localization maps send $F/X_i^d$ and $F/X_k^d$ to fractions
in the SAME localization of the product $X_iX_k$.  Equality after cross
multiplication is the literal polynomial identity; the native localization
equivalence proves equality even with zero divisors and nilpotents.
The actual homogeneous away-to-section map gives the whole scheme section.
Its native restriction naturality transports the preceding equality to the
whole original overlap.  In degree one, the original hyperplane local-equation
formula is exactly this same homogeneous fraction, so the unit in the overlap
law is the actual original coordinate transition unit.
\end{proof}

\begin{theorem}[Every original homogeneous form constructs a global actual O(d) section]
\label{thm:fc-original-homogeneous-form-section}
\leanok
\lean{EnvelopeProperty.ProjectiveHomogeneousSections.localSection, EnvelopeProperty.ProjectiveHomogeneousSections.localSection_overlap, EnvelopeProperty.ProjectiveHomogeneousSections.exists_section, EnvelopeProperty.ProjectiveHomogeneousSections.homogeneousSection, EnvelopeProperty.ProjectiveHomogeneousSections.section_restrict}
\uses{lem:fc-original-homogeneous-form-coefficients}
Every original homogeneous polynomial $F$ of degree $d\geq0$ constructs a
genuine global section of the actual native coordinate-hyperplane twist
$\mathcal O(d)$.  On EVERY whole original coordinate chart its restriction
is precisely $F/X_i^d$ times the actual original coefficient-one line frame.
No line-section certificate, finiteness of the coordinate index, or domain
hypothesis is supplied.
\end{theorem}
\begin{proof}
Multiply the preceding actual coefficient section by the actual original
coefficient-one frame of the native positive twist on each original chart.
The genuine original frame-transition theorem says that $(X_k/X_i)^d$
times the restricted $i$-frame is the restricted $k$-frame.  Combining it
with the preceding original scalar-transition law makes the whole local
line sections equal on the whole original overlap.  The actual coordinate
opens cover native projective space.  Genuine module-sheaf gluing therefore
constructs a global actual section with exactly those original restrictions.
The native gluing restriction identity gives the stated formula for every chart.
\end{proof}

\begin{lemma}[The original homogeneous section has its exact original line-frame coefficient]
\label{lem:fc-original-homogeneous-form-frame-coefficient}
\leanok
\lean{EnvelopeProperty.ProjectiveHomogeneousSections.section_coefficient, EnvelopeProperty.ProjectiveHomogeneousSections.sectionHom}
\uses{thm:fc-original-homogeneous-form-section}
The actual original line-frame coefficient of this global section on the
original $i$-chart is exactly the actual original scheme fraction $F/X_i^d$.
The same global section defines a genuine whole morphism from the original
structure line to the actual original $\mathcal O(d)$.
\end{lemma}
\begin{proof}
Use the exact original restriction formula of the preceding theorem.
The genuine inverse local frame sends $1$ to the coefficient-one frame;
applying the actual frame to its scalar multiple returns the same original
scalar by module linearity and the actual inverse identities.  Thus the
coefficient is precisely $F/X_i^d$, without an identification assumption.
The original module global-section/unit-Hom equivalence gives the genuine
whole structure-line morphism attached to this same original global section.
\end{proof}
